\section{Background and the Symmetry Collapse}
\label{sec:background}

\subsection{Standard ReCom}

Let $G = (V, E)$ be the census-tract adjacency graph for a state, and let
$\pi : V \to [k]$ be a valid redistricting plan (each district connected
and population-balanced within tolerance $\epsilon$).
A single \recom{} step proceeds as follows~\citep{deford2021}:

\begin{enumerate}
  \item Sample an adjacent pair of districts $(d_i, d_j)$ uniformly at
        random from all pairs sharing at least one boundary edge.
  \item Form the merged region $R = \pi^{-1}(d_i) \cup \pi^{-1}(d_j)$,
        with the induced subgraph $G[R]$.
  \item Sample a uniformly random spanning tree $T$ of $G[R]$ via
        Wilson's algorithm~\citep{wilson1996}.
  \item Enumerate the balanced cuts of $T$ --- edges whose removal partitions
        $T$ into two components, each satisfying the population balance
        constraint.
  \item If at least one balanced cut exists, sample one uniformly and
        accept the resulting plan $\pi'$.
        Otherwise, reject and repeat from step~1 (pair reselection).
\end{enumerate}

Step~3 is the dominant cost: Wilson's algorithm runs in $O(|R| \log |R|)$
expected time for a graph with $|R|$ vertices.

\subsection{Why the Acceptance Ratio Is Always 1}

The Metropolis-Hastings acceptance ratio for the move $\pi \to \pi'$ is:
\[
  \alpha(\pi, \pi') = \min\!\left(1,\;
    \frac{\pi^*(\pi')\, Q(\pi', \pi)}{\pi^*(\pi)\, Q(\pi, \pi')}
  \right),
\]
where $\pi^*$ is the target distribution and $Q(\pi, \pi')$ is the proposal
probability.
If we target the uniform distribution $\pi^*(\pi) = \pi^*(\pi')$,
the ratio simplifies to $Q(\pi', \pi) / Q(\pi, \pi')$.

The proposal probability for standard \recom{} is:
\begin{align}
  Q(\pi, \pi') &= \frac{1}{|\mathcal{P}(\pi)|}
    \cdot \frac{1}{|\mathcal{T}(R)|}
    \cdot \frac{1}{|\mathcal{C}(T, R)|},
\end{align}
where $\mathcal{P}(\pi)$ is the set of adjacent pairs of districts in $\pi$,
$\mathcal{T}(R)$ is the set of spanning trees of $G[R]$, and
$\mathcal{C}(T, R)$ is the set of balanced cuts of the sampled tree $T$.
The reverse proposal $Q(\pi', \pi)$ has the same structure, with the
same merged region $R$ (since $\pi$ and $\pi'$ differ only in how
$R$ is partitioned).

In standard \recom{}, both the forward and reverse proposals sample a spanning
tree from the same graph $G[R]$ with the same uniform measure.
Consequently:
\[
  \frac{Q(\pi', \pi)}{Q(\pi, \pi')}
  = \frac{|\mathcal{P}(\pi)| \cdot |\mathcal{C}(T', R)|}{|\mathcal{P}(\pi')|
    \cdot |\mathcal{C}(T, R)|}.
\]
Here $T$ is the spanning tree used to propose $\pi'$, and $T'$ is the
(hypothetical) spanning tree that would propose $\pi$ from $\pi'$.
Standard \recom{} does \emph{not} sample $T'$---it only ever samples $T$.
Without an independent sample of $T'$, the ratio $|\mathcal{C}(T', R)|$ is
unknown, and the ratio $Q(\pi', \pi) / Q(\pi, \pi')$ is not computed.
The implementation effectively sets $\alpha = 1$ (always accept if a balanced
cut exists), which is equivalent to using an unknown acceptance ratio.

The stationary distribution of the resulting chain is not the uniform
distribution; it is a distribution that weights plans by the implicit bias
of the spanning tree measure and the number of balanced cuts of the typical
tree in the region.

\subsection{Breaking the Symmetry via Independent Reverse Tree}

Merge-Split~\citep{janson2022} restores computable Metropolis-Hastings
ratios by sampling $T'$---the reverse spanning tree---independently:

\begin{enumerate}
  \item After sampling $T_{\mathrm{fwd}}$ and selecting the proposed plan
        $\pi'$, sample a second uniformly random spanning tree
        $T_{\mathrm{rev}}$ of the same merged region $G[R]$, using an
        independent random source.
  \item Count $\mathrm{cuts\_fwd} = |\mathcal{C}(T_{\mathrm{fwd}}, R)|$
        and $\mathrm{cuts\_rev} = |\mathcal{C}(T_{\mathrm{rev}}, R)|$.
  \item Use the ratio $\mathrm{cuts\_fwd} / \mathrm{cuts\_rev}$ as the
        acceptance probability.
\end{enumerate}

The key insight is that $|\mathcal{C}(T_{\mathrm{rev}}, R)|$ is a random
estimate of the expected balanced-cut count $\mathbb{E}[|\mathcal{C}(T', R)|]$
over all spanning trees $T'$.
Since both trees are sampled uniformly, the ratio is an unbiased estimator
of the true ratio $Q(\pi', \pi) / Q(\pi, \pi')$ in expectation, giving
approximate detailed balance.

\begin{remark}[Exact vs.\ approximate reversibility]
\label{rem:approx}
The \mergesplit{} ratio is a random estimate, not the exact matrix-tree ratio
used by \forestrecom{}~\citep{mccartan2020,dellaLibera2026forestRecom}.
The approximation is \emph{consistent}---the chain still has a stationary
distribution---but the stationary distribution is the uniform distribution
only in expectation over the randomness of $T_{\mathrm{rev}}$.
For practical redistricting ensembles this distinction rarely matters
(both produce similar distributional profiles in our experiments), but
it is important for theoretical claims about exact uniformity.
\end{remark}
